\section{Extensibility}
\label{sec:5}
This section explores two extensions on top of canonical retrieval models, and discuss how we scale out SilverTorch to multi-GPUs.

\subsection{OverArch Scoring}
In the two-tower model architecture, user-item similarities are computed using dot product. Unfortunately, this approach has no trainable parameters and oversimplifies user-item interactions. Without SilverTorch's GPU-based approach, meeting latency requirements while adding complex interaction layers is challenging. SilverTorch substantially reduces inference cost, enabling retrieval models to incorporate scoring layers beyond nearest neighbor search.
Building on the ANN search and feature filtering queries defined in Section~\ref{sec:4}, a SilverTorch model includes additional scoring layers (referred to as the OverArch layer). The retrieval process operates in two steps. First, it executes the query to pre-filter $K_0$ items (on the order of $O(10{,}000)$ to $O(100{,}000)$) from original candidate pool, where dot product is still adopted for distance calculation within ANN search. Second, the OverArch employs neural network modules to rank the $K_0$ user-item pairs and returns the final retrieval results. We find that introducing OverArch achieves better recall for retrieval than enhancing ANN search accuracy alone.
An OverArch layer can be defined as a Multi-layer Perceptron (MLP), or a multiple stacked self-attention layers to capture the correlation to understand user’s interest, with the capability of looking at items in an entire session. Recent study proposes more structured interaction layer using Mixture of logits (MoL)~\cite{zhai2023revisiting}~\cite{ding2025retrieval} that defines similarity as adaptive composition of elementary functions. Benefiting from SilverTorch's caching design, item embeddings and cross-features used in the OverArch can be directly extracted from GPU memory during publish. SilverTorch supports complex OverArch architectures that could improve model quality of retrieval through the full fidelity of training and serving consistency.

\subsection{Multi-Task Retrieval with Value Model}
To learn multiple aspects of users, recommendation should predict multiple objectives to capture diverse user behaviors such as content—like, share, or comment. Multi-task learning addresses this by training a unified model that exploits commonalities and differences across tasks, improving prediction accuracy through knowledge sharing. However, previous retrieval systems largely avoided multi-task approaches due to prohibitive latency costs.

A straightforward idea is to apply different user embeddings and item embeddings to represent each task. However, item embeddings are learned to be the semantics representation so they should be shared across tasks. Therefore, in multi-task retrieval, the user tower shares a lookup table but applies task-specific dense layers to generate user embeddings per task, while item embeddings are shared across tasks.
At serving time, ANN search handles multi-task queries and returns task-specific item lists. CPU-based solutions require replicating the ANN index for each task to avoid latency increases, causing linear cost growth. In contrast, SilverTorch leverages GPU parallelism to batch requests from multiple tasks within a single index copy without latency regression, making multi-task retrieval cost-efficient.
After ANN search and filtering, a merge operation combines results across tasks before the OverArch layer predicts per-task scores. To aggregate these predictions into a unified engagement score, we introduce a Value Model (VM) that transforms the model predictions to business objectives. The VM combines predictions (likes, shares, comments) through user-defined formulas expressed in a JSON-like format with pre-defined conditions. SilverTorch implements a GPU VM kernel that parses formulas into an abstract syntax tree and process multiple items in parallel. Applying value-based aggregation at retrieval provides better consistency between retrieval and ranking.

\subsection{Scale Out}
To handle larger candidate pools and models, SilverTorch scales out to multiple GPUs. The ANN index and Bloom index are sharded across GPU cards, with each GPU processing a partition of items, while OverArch parameters are replicated on each GPU. During serving, each GPU independently computes local pre-filtered results through ANN search and feature filtering. Item embeddings are then gathered to a single GPU to compute the final retrieval results. Since each GPU maintains a copy of the OverArch and Value Model, request batches can be evenly distributed across GPUs in parallel, preventing any single GPU from becoming a bottleneck.

GPU memory is the primary factor determining the number of GPUs required for serving. It depends on multiple factors: ANN and Bloom index size and weights of the OverArch. As candidate pool grows, both of Bloom and ANN index scale accordingly. ANN index size further depends on quantization precision and embedding dimensions. SilverTorch's Int8 ANN and compact Bloom index designs significantly reduce GPU memory footprint. The complexity of the OverArch also constrains how many items a single GPU can serve before reaching limits. For user embedding tables, we leverage CPU-based parameter servers for distributed inference.
Adopting the unified design facilitates co-design across recommendation services, enabling future flexible disaggregation of serving components based on serving characteristics (e.g., compute-bound vs. memory-bound) rather than pre-defined boundaries. Additionally, real traffic is bursty, requiring automatic GPU scaling based on offline capacity estimation. We support QPS-based scaling that automatically scales GPU instances up or down within minutes. For extreme bursts exceeding capacity, excess traffic is throttled.